\documentclass[journal]{IEEEtran}
\usepackage{amsmath,amssymb,amsfonts}
\usepackage{graphicx}
\usepackage{booktabs}
\usepackage{hyperref}
\usepackage{float}

\begin{document}

\title{Supplementary Material for ResQ-AI: An Uncertainty-Aware Multimodal Framework}

\author{ResQ-AI Team}

\maketitle

\section{Extended Hyperparameter Configurations}
\label{sec:hyperparams}
This section details the hyperparameter search spaces and final selected values for the ResQNet model, UQ methods, and the stochastic optimization module.

\begin{table}[H]
\centering
\caption{Hyperparameters for ResQNet Training}
\begin{tabular}{llc}
\toprule
Parameter & Search Space & Selected \\
\midrule
Learning Rate & [1e-5, 1e-4, 1e-3] & 3e-4 \\
Batch Size & [8, 16, 32] & 16 \\
Weight Decay & [1e-6, 1e-5, 1e-4] & 1e-5 \\
Modality Dropout Prob & [0.0, 0.2, 0.5] & 0.2 \\
\bottomrule
\end{tabular}
\end{table}

\section{Mathematical Derivations}
\label{sec:math}
\subsection{Chance Constraint Reformulation}
The chance constraint for satisfying demand $P(\sum x_{mn} \ge D_n) \ge 1 - \alpha$ is reformulated using Big-M constraints in the Sample Average Approximation framework.
Let $y^{(s)} \in \{0,1\}$ indicate if the constraint is violated in scenario $s$.
$\sum_{m} x_{mn} \ge D_n^{(s)} - M y^{(s)}$ and $\sum_s y^{(s)} \le \alpha S$.

\section{Additional Qualitative Results}
\label{sec:qualitative}
Additional visualizations of the uncertainty maps are provided here.

\end{document}
